%----------------------------------------------------------------------------- %%
\section{Literature Review}
\label{bab:2}
%----------------------------------------------------------------------------- %%

This section reviews the theoretical concepts and prior work that form the foundation of this study, including Kubernetes resource management, resource fragmentation, Descheduler mechanisms, and predicted-target eviction-candidate selection.

\subsection{Kubernetes Architecture and Resource Management}
\label{sec:kubernetesArchitecture}

Kubernetes distributes Pods across worker Nodes. For CPU and memory, a container may declare a resource request, a resource limit, both, or neither. Requests represent the resource baseline used by the scheduler when determining whether a Pod fits on a Node, whereas limits are runtime enforcement ceilings applied by the kubelet and container runtime~\cite{k8s_resource_management}. Kubernetes uses effective Pod requests when evaluating placement feasibility and resource allocation on each Node. Resource limits do not participate in this fit calculation; they define runtime enforcement ceilings after the Pod is running. Scheduler scoring strategies such as \textit{MostAllocated} and \textit{RequestedToCapacityRatio} may then use request-based allocation to favor denser packing~\cite{k8s_resource_bin_packing}.

Requests and limits also contribute to Pod Quality of Service (QoS) classification. A Pod is \textit{Guaranteed} when every container declares CPU and memory requests equal to its corresponding limits, while Pods with partial or unequal declarations are generally \textit{Burstable}. A Pod whose containers declare neither CPU nor memory requests or limits is \textit{BestEffort}~\cite{k8s_pod_qos}. If a limit is declared without a request, Kubernetes can use that limit as the effective request for the resource. A namespace \textit{LimitRange} may also assign default requests or limits during Pod admission. These defaults become part of the Pod's effective resource configuration~\cite{k8s_limit_range}.

These decisions are made when a Pod is scheduled. A running Pod is not continuously relocated when the cluster state changes, so an initially acceptable layout can become inefficient after scale-outs, terminations, and replacement events. This separation between initial placement and later cluster state motivates post-placement correction.

\subsection{Resource Fragmentation and Stranded Capacity}
\label{sec:resourceFragmentation}

Resource fragmentation occurs when free capacity is divided among Nodes in shapes that do not match incoming workloads. A cluster can have enough aggregate CPU and memory while no individual Node satisfies both dimensions of a Pod request. This makes CPU and memory placement a multi-dimensional problem rather than a scalar load-balancing problem.

CPU-memory skew is one source of stranded capacity. If one Node is CPU-heavy and another is memory-heavy, each leaves a different residual resource that cannot be consumed alone. The absolute difference between normalized CPU and memory fractions provides a direct shape indicator, while the product of free CPU and memory fractions represents the jointly usable free-space region. Prior work shows that directed eviction can repair fragmented Kubernetes layouts and support Node reclamation~\cite{larsson_klein_elmroth_2023}.

\subsection{Resource Defragmentation and Multi-Resource Bin Packing}
\label{sec:binPacking}

Resource packing seeks to place workloads on fewer active Nodes so that empty Nodes can be removed or powered down. In a multi-resource setting, dense packing alone is insufficient: a target can be highly utilized but badly skewed. A useful packing score should therefore represent both density and CPU-memory balance. Container resource defragmentation and task scheduling have been studied using optimization and learning-based methods~\cite{guo2025joint,bagaa2024multiobjective}, but these are commonly oriented toward placement or scheduling decisions rather than long-term post-placement correction in Kubernetes clusters.

The policy in this research uses normalized CPU and memory fractions. Average utilization represents density, while a bin score combines density with the absolute distance between the two dimensions. The Resource Imbalance Index (RII) preserves the direction of skew, and the Free-Space Index (FSI) measures the rectangular CPU-memory headroom. Together these quantities identify Nodes that are lightly loaded, poorly shaped, or under pressure from a narrow free-space region.

\subsection{Kubernetes Descheduler}
\label{sec:deschedulerLimitations}

The Kubernetes Descheduler evicts Pods that violate a policy or contribute to an undesirable placement state, after which normal controllers recreate the Pods and kube-scheduler places them again~\cite{k8s_descheduler_docs}. It provides an appropriate extension point for node reclamation because it can repair placement without modifying the scheduler.

An eviction policy must nevertheless account for the scheduler boundary. Deleting a Pod does not guarantee that it moves to a useful Node; it may become Pending, return to its origin, or land on a target that worsens fragmentation. A node reclamation-oriented policy therefore benefits from predicting feasible non-origin targets and rejecting movements that exceed a configured ceiling or pack downward.

Upstream utilization policies provide important baselines. \textit{HighNodeUtilization} packs by draining low-utilization Nodes toward more utilized Nodes, but it does not explicitly optimize CPU-memory stranding. This distinction is central to the evaluation in Section~\ref{bab:4}.

\subsection{Predicted-Target Eviction-Candidate Selection}
\label{sec:predictedTargetSelection}

Eviction-candidate selection can be formulated around the expected landing state. For every evictable Pod on a drain candidate, \textit{predictSchedulerTarget} simulates request-based placement on non-origin Nodes that satisfy scheduler constraints. Feasible targets must have sufficient CPU and memory, remain below a node reclamation ceiling after placement, and be at least as packed as the source.

The target with the best projected CPU-memory bin score is retained for each Pod. The Pod whose predicted target has the highest score is selected for eviction. This single criterion aligns the action directly with the desired destination state: the selected movement should place complementary resource shapes together while preserving node reclamation pressure.

\subsection{Current Approaches and Research Gap}
\label{sec:researchGap}

Existing mechanisms address individual parts of the problem. Scheduler scoring improves new placement, utilization-based descheduling drains Nodes according to thresholds, and fragmentation research motivates directed eviction. The remaining gap is a practical Descheduler policy that combines:

\begin{itemize}
    \item source assessment using utilization and CPU-memory shape;
    \item RII/FSI-based drain ordering;
    \item request-feasible scheduler-target prediction;
    \item a direct projected-target balance criterion; and
    \item bounded, convergent eviction.
\end{itemize}

\textit{ResourceDefragmentation} is designed to fill this gap while leaving final binding to the standard kube-scheduler.
